
\section{一系列关键工具}

现在，我们已经准备好以一种动手的方式来理解\trCurve{}了。我们将反复使用下面这一系列关键评注；而把它们挂在嘴边、随时能信手拈来，是很要紧的。

\textbf{19.2.1.}\quad \emph{在下文中，$C$ 将是域 $k$ 上一条\trProjective{}、几何正则、\trGeometricallyIntegral{}的\trCurve，而 $\mathscr{L}$ 是 $C$ 上的一个\trInvertibleSheaf。}

\textbf{19.2.2. 提醒：\trSerreDuality。}\quad $k$ 上\trGeometricallyIrreducible{}正则的\trGenus{} $g$ \trCurve{} $C$ 上的\trSerreDuality（\mbox{定理~18.5.1}）牵涉到一个\trInvertibleSheaf{} $\omega_C$（次数为 $2g-2$，且带有 $g$ 个线性无关的截面，\mbox{练习~18.5.A}），使得对 $C$ 上任意有限\trRank{}\trLocallyFreeSheaf{} $\mathscr{F}$，当 $i=0$、$1$ 时都有 $h^i(C,\mathscr{F})=h^{1-i}(C,\omega_C\otimes\mathscr{F}^{\vee})$。（更好的说法：在那两个\trCohomology{}群之间存在一个对偶。）

\textbf{19.2.3. 负\trDegree{}\trLineBundle{}没有非零\trSection。}\quad 若 $\deg\mathscr{L}<0$，则 $h^0(C,\mathscr{L})=0$。理由：$\deg\mathscr{L}$ 是任一\trRational{}截面的零点个数减去极点个数（适当地计数）（\mbox{重要练习~18.4.C}）。若存在一个正则截面（即没有极点的截面），此时该次数必非负。把这番论证加以精化，便得到下面的结论。

\textbf{19.2.4. \trDegree{} $0$ 的\trLineBundle，以及识别它们何时是\trTrivial{}的。}\quad 若 $\deg\mathscr{L}=0$，则 $h^0(C,\mathscr{L})=0$ 或 $1$；而且若 $h^0(C,\mathscr{L})=1$，则 $\mathscr{L}\cong\OO_C$。理由：若存在一个非零截面 $s$，则它没有极点，从而也没有零点，因为 $\deg\mathscr{L}=0$。于是 $\operatorname{div}s=0$，故 $\mathscr{L}\cong\OO_C(\operatorname{div}s)=\OO_C$。（回想一下这是怎么回事；参看\mbox{重要练习~15.4.G}：$s$ 给出该\trInvertibleSheaf{}的一个\trTrivialization。对任意开集 $\mathrm{U}$\footnote{原书此处作「for any open set」，当为「for any open set U」之脱文，径补。——译者注}，我们都有自然的双射 $\Gamma(\mathrm{U},\mathscr{L})\overset{\sim}{\longleftrightarrow}\Gamma(\mathrm{U},\OO_{\mathrm{U}})$，其中从左到右的映射是 $s'\mapsto s'/s$，从右到左的映射是 $f\mapsto sf$。）反之，对几何整的射影\trVariety{}而言，$h^0(\OO)=1$。（\S11.4.7 在 $k$ \trAlgClosed{}时证明了这一点——这正是用到几何整性的地方——而\mbox{练习~18.2.H} 则证明了\trCohomology{}与基域扩张交换。）

\trSerreDuality{}把上述关于\trDegree{}至多为 $0$ 的\trLineBundle{}的诸陈述，转成了关于\trDegree{}至少为 $2g-2$ 的\trLineBundle{}的诸陈述，如下。

\textbf{19.2.5. 若\trDegree{}足够高，则我们知道 $h^0(C,\mathscr{L})$。}\quad 若 $\deg\mathscr{L}>2g-2$，则
\begin{equation*}\tag{19.2.5.1}\label{eq:h0-high-degree}
\fbox{$h^0(C,\mathscr{L})=\deg\mathscr{L}-g+1$.}
\end{equation*}
于是当 $\deg\mathscr{L}\gg 0$ 时，我们就知道 $h^0(C,\mathscr{L})$。（\emph{这很重要——要记住这一条！}）理由：$h^1(C,\mathscr{L})=h^0(C,\omega_C\otimes\mathscr{L}^{\vee})$；但 $\omega_C\otimes\mathscr{L}^{\vee}$ 的\trDegree{}为负（因为 $\omega_C$ 的\trDegree{}是 $2g-2$），于是这个\trInvertibleSheaf{}没有截面。这个结果随即由以 \eqref{eq:riemann-roch-h0-h1} 的面目出现的\trRiemannRoch{}定理得到。

\textbf{19.2.A.（有用练习：在\trDegree{} $2g-2$ 的诸\trLineBundle{}中认出 $\omega_C$）}\quad 设 $\mathscr{L}$ 是一个\trDegree{}为 $2g-2$ 的\trInvertibleSheaf。证明：它有 $g-1$ 或 $g$ 个截面，而且它有 $g$ 个截面，当且仅当 $\mathscr{L}\cong\omega_C$。

\textbf{19.2.B.（练习）}\quad 设 $C$ 是一条\trGenus{} $g>1$ 的\trCurve，所在的域 $k$ 不必\trAlgClosed。证明：$C$ 有一个在基域上\trDegree{}至多为 $2g-2$ 的\trClosedPoint。（作为对照：当 $g=1$ 时，对任意 $n$，都存在 $\mathbb{Q}$ 上的一条\trGenus{} $1$ \trCurve，它没有\trDegree{}小于 $n$ 的点，[MO286390]\footnote{译注：[MO286390] 指 MathOverflow 问答第 286390 号。}！）

\textbf{19.2.6. 用一个（\trDegree{} $1$ 的）点扭转 $\mathscr{L}$，至多把 $h^0$ 改变 $1$。}\quad 设 $p$ 是任意一个\trDegree{}为 $1$ 的\trClosedPoint（即 $p$ 的\trResidueField{}是 $k$）。则 $h^0(C,\mathscr{L})-h^0(C,\mathscr{L}(-p))=0$ 或 $1$。（用一个\trDivisor{}扭转 $\mathscr{L}$，比如 $\mathscr{L}(-p)$，曾定义于 \S15.4.17。）理由：考察 $0\to\OO_C(-p)\to\OO_C\to\OO|_{p}\to 0$；与 $\mathscr{L}$ 作张量（由于 $\mathscr{L}$ \trLocallyFree，这是正合的），便得到
\[
0\longrightarrow \mathscr{L}(-p)\longrightarrow \mathscr{L}\longrightarrow \mathscr{L}|_{p}\longrightarrow 0.
\]
那么 $h^0(C,\mathscr{L}|_{p})=1$，于是\trCohomology{}的\trLongExactSequence{}开头为
\[
0\longrightarrow \mathrm{H}^0(C,\mathscr{L}(-p))\longrightarrow \mathrm{H}^0(C,\mathscr{L})\longrightarrow \mathrm{H}^0(C,\mathscr{L}|_{p}),
\]
我们便大功告成。

\textbf{19.2.7. $\mathscr{L}$ 无基点的判别法。}\quad 就本条评注而言，设 $k$ \trAlgClosed，从而\emph{一切}\trClosedPoint{}在 $k$ 上的\trDegree{}都是 $1$。那么，若对\emph{一切}\trClosedPoint{} $p$ 都有 $h^0(C,\mathscr{L})-h^0(C,\mathscr{L}(-p))=1$，则 $\mathscr{L}$ \trBasePointFree，从而诱导出一个从 $C$ 到\trProjectiveSpace{}的\trMorphism（\mbox{定理~15.2.2}）。理由：任给 $p$，我们的等式表明，存在 $\mathscr{L}$ 的一个在 $p$ 处不消失的截面——于是按定义，$p$ 不是 $\mathscr{L}$ 的\trBasePoint。

\textbf{19.2.8.}\quad 接着，设 $p$ 与 $q$ 是互不相同的两个（闭）点，\trDegree{}皆为 $1$。则 $h^0(C,\mathscr{L})-h^0(C,\mathscr{L}(-p-q))=0$、$1$ 或 $2$（把\mbox{注记~19.2.6} 的论证重复两遍即可）。若 $h^0(C,\mathscr{L})-h^0(C,\mathscr{L}(-p-q))=2$，则必然
\begin{equation*}\tag{19.2.8.1}\label{eq:separating-points-equality}
h^0(C,\mathscr{L})=h^0(C,\mathscr{L}(-p))+1=h^0(C,\mathscr{L}(-q))+1=h^0(C,\mathscr{L}(-p-q))+2.
\end{equation*}
这时，线性系 $\mathscr{L}$ 分离点 $p$ 与 $q$，也就是说，到\trProjectiveSpace{}的相应映射 $f$ 满足 $f(p)\ne f(q)$。理由：\trProjectiveSpace{}中存在一张\trHyperplane{}经过 $f(q)$ 而不经过 $f(p)$；等价地，$\mathscr{L}$ 存在一个在 $q$ 处消失而在 $p$ 处不消失的截面。这正是 \eqref{eq:separating-points-equality} 中最后一个等式的内容。

\textbf{19.2.9.}\quad 由与上文相同的论证，若 $p$ 是一个\trDegree{}为 $1$ 的（闭）点，则 $h^0(C,\mathscr{L})-h^0(C,\mathscr{L}(-2p))=0$、$1$ 或 $2$。我宣称：若这个值是 $2$，则相应于 $\mathscr{L}$ 的映射（由上文已知它在 $p$ 处无基点）分离 $p$ 处的切向量。为了证明这一点，我们需要证明余切映射是满射。而要证明到一个一维\trVectorSpace{}的满射性，只须证明该映射非零。于是，只须在靶上找到一个在 $p$ 的像处消失的函数，使得其拉回是一个在 $p$ 处消失到 $1$ 阶而非 $2$ 阶的函数。换言之，我们想要的是 $\mathscr{L}$ 的一个在 $p$ 处消失到 $1$ 阶而不消失到 $2$ 阶的截面。而这正是陈述 $h^0(C,\mathscr{L}(-p))-h^0(C,\mathscr{L}(-2p))=1$ 的内容。

\textbf{19.2.10. $\mathscr{L}$ 为极丰富的判别法。}\quad 把上文的一些评注结合起来：设 $C$ 是代数闭域 $k$ 上的一条\trCurve，而 $\mathscr{L}$ 是一个\trInvertibleSheaf，使得对\emph{一切}\trClosedPoint{} $p$ 与 $q$——\emph{不必互不相同}——都有 $h^0(C,\mathscr{L})-h^0(C,\mathscr{L}(-p-q))=2$；那么 $\mathscr{L}$ \emph{给出一个到\trProjectiveSpace{}的\trClosedEmbedding}，因为它分离点与切向量：由\mbox{定理~19.1.1}。

\textbf{19.2.C.（练习）}\quad 假设 $k$ \trAlgClosed，于是前一条评注适用。证明：$C\setminus\{p\}$ 是\trAffine{}的。提示：证明若 $j\gg 0$，则 $\OO(jp)$ 无基点，并且至少有两个线性无关的截面，其中之一的\trDivisor{}是 $jp$。用这两个截面映射到 $\mathbb{P}^1$，使得 $\infty$ 的集合论原像是 $p$。论证该映射是\trFinite{}的，而 $C\setminus\{p\}$ 是 $\mathbb{A}^1$ 的原像。（这一论证的一个平凡变体表明：当 $n>0$ 时，$C\setminus\{p_1,\dots,p_n\}$ 是\trAffine{}的。）

\textbf{19.2.11. 有力结论：$\mathscr{L}$ 无基点与极丰富的数值判别法。}\quad 我们可以把上述讨论的很大一部分合并起来，给出下面这条有用的事实。若 $k$ \trAlgClosed，则 $\deg\mathscr{L}\ge 2g$ 蕴涵 $\mathscr{L}$ \trBasePointFree（从而确定一个到\trProjectiveSpace{}的\trMorphism）。此外，$\deg\mathscr{L}\ge 2g+1$ 蕴涵：对 $\mathscr{L}$ 的任意一组截面基 $s_0,\dots,s_{\deg\mathscr{L}-g}$，由 $[s_0,\dots,s_{\deg\mathscr{L}-g}]$ 给出的\trMorphism{}事实上是一个到 $\mathbb{P}^{\deg\mathscr{L}-g}$ 的\trClosedEmbedding（于是 $\mathscr{L}$ 是\trVeryAmple{}的）。记住这一条！

\textbf{19.2.D.（练习：关于 $\bar{k}$ 上射影、正则、整的曲线，\trAmple{}＝正\trDegree）}\quad 证明：$\bar{k}$ 上一条射影、正则、整的曲线上的\trInvertibleSheaf{} $\mathscr{L}$ 是\trAmple{}的，当且仅当 $\deg\mathscr{L}>0$。

（借下文\mbox{练习~19.2.E} 之助，这一点可以推广到一般域上的曲线。）于是，曲线上\trLineBundle{}的\trAmpleness{}有一条简单粗暴的纯数值判别法。这一点可以推广到更高维数的射影\trVariety；它称作\emph{丰富性的\trNakai--\trMoishezon{}判别法}（Nakai--Moishezon criterion for ampleness），见\mbox{定理~20.4.1}。

\textbf{19.2.E.（练习：向非\trAlgClosed{}域的扩张）}\quad 证明：即使去掉 $k$ \trAlgClosed{}这一假设，\S19.2.11 中的诸陈述也依然成立。（提示：要证明关于某条\trCurve{} $C$ 与某个\trLineBundle{} $\mathscr{L}$ 的两条事实中的一条，改为考察 $C\times_{\Spec k}\Spec\bar{k}$。然后证明：若 $\mathscr{L}$ 在这里的拉回具有能给出两条所欲性质之一的截面，则「底下」就有具有同样性质的截面。你或许想动用我们已经用过的事实，例如：无基点性独立于基域的任何扩张（\mbox{练习~18.2.I}）；或者，$k$ 上的仿射态射为\trClosedEmbedding{}这一性质，在作 $k$ 的扩张之后成立，当且仅当它在扩张之前成立（\mbox{练习~10.2.K}）。）

\textbf{19.2.F.（练习：扩充\mbox{练习~19.2.D}）}\quad 设 $\mathscr{L}$ 是（域 $k$ 上的）一条\trProjective{}、几何正则、\trGeometricallyIntegral{}\trCurve{} $C$ 上的一个\trInvertibleSheaf。证明：$\mathscr{L}$ 是\trAmple{}的，当且仅当 $\deg\mathscr{L}>0$。提示：借\mbox{练习~10.2.K} 之助，化归到 $k$ \trAlgClosed{}的情形。

\textbf{19.2.G.（练习）}\quad 设 $C$ 是一条\trCurve（照本节的惯例，即取 \S19.2.1 意义下的曲线），而 $p\in C$ 是一个\trClosedPoint。证明：$C\setminus\{p\}$ 是\trAffine{}的。提示：\mbox{练习~19.2.F}。

现在，我们已整装待发，要带上这些事实到 \S19.3 中大干一场了。

\textbf{19.2.12.}\ \hard\ \textbf{\trhRemark：\trProper{}\trCurve{}是\trProjective{}的。}

（在 \S19.2 的这最后一部分中，我们不再假设诸曲线是光滑的，甚至不再假设它们正则或既约。）我们此处的讨论，连同\trSerre{}的丰富性上同调判别法（\S18.8），是证明「域上每条\trProper{}\trCurve{}都是\trProjective{}的」这一结果的关键配料。

\textbf{19.2.H.（练习）}\quad 证明：每条\trReduced{}的\trProper{}\trCurve{} $C$ 都是\trProjective{}的。提示：在 $C$ 的每个\trIrreducibleComponent{}上选一个正则点，并令 $\mathscr{L}$ 为相应的\trInvertibleSheaf。我们盼着 $\mathscr{L}$ 会是\trAmple{}的。回想：$C$ 上的一个\trLineBundle{}是\trAmple{}的，若它到 $C$ 的\trNormalization{}上的拉回是\trAmple{}的（\mbox{棘手练习~18.8.C}）。所以请证明：到\trNormalization{}上的拉回是\trAmple{}的。

\textbf{19.2.I.（练习）}\quad 先在 $C$ 上找出一个你将证明为\trAmple{}的\trInvertibleSheaf，然后利用\mbox{练习~18.8.A} 证明：每条\trProper{}（不必\trReduced）的\trCurve{} $C$ 都是\trProjective{}的。

\textbf{19.2.J.（练习）}\quad 证明：射影曲线上的一个\trLineBundle{}是\trAmple{}的，当且仅当它在每个分量上的\trDegree{}都是正的。
